% Part: methods
% Chapter: proofs
% Section: resources

\documentclass[../../../include/open-logic-section]{subfiles}

\begin{document}

\olfileid{mth}{prf}{res}

\section{અન્ય સ્રોતો}

ગણિતમાં સાબિતીઓ કેવી રીતે કરવી તે અંગે ઘણાં ઉપયોગી પુસ્તકો
છે. ખાસ કરીને \emph{How to Read and do Proofs: An Introduction to
Mathematical Thought Processes} \citep{Solow2013} અને \emph{How to
Prove It: A Structured Approach} \citep{Velleman2019} જુઓ.
\href{http://www.people.vcu.edu/~rhammack/BookOfProof/BookOfProof.pdf}{\emph{Book
of Proof}} \citep{Hammack2013} અને
\href{https://scholarworks.gvsu.edu/books/7/}{\emph{Mathematical
Reasoning}} \citep{Sandstrum2019} એ સાબિતી અંગેનાં પુસ્તકો
ઓનલાઇન નિઃશુલ્ક ઉપલબ્ધ છે. તત્ત્વચિંતકોને ગાણિતિક તર્કવિચારના
પ્રારંભિક પરિચય માટે
\emph{More
Precisely: The Math you need to do Philosophy} \citep{Steinhart2018}
ઉપયોગી લાગી શકે.

ઇન્ટરનેટ પર સાબિતી અંગે ટૂંકી માર્ગદર્શિકાઓ પણ છે;
દાખલા તરીકે,
\href{https://math.berkeley.edu/~hutching/teach/proofs.pdf}{``Introduction
  to Mathematical Arguments''} \citep{Hutchings2003} અને
\href{https://eugeniacheng.com/wp-content/uploads/2017/02/cheng-proofguide.pdf}{``How to
  write proofs''} \citep{Cheng2004}.

\subsection{પ્રેરક વિડિયો}

ઘરકામ કરવા ઉત્સાહ નથી લાગતો? મન ઉદાસ છે? કદાચ આ વિડિયો
મદદરૂપ થાય!

\begin{itemize}
\item \url{https://www.youtube.com/watch?v=ZXsQAXx_ao0}
\item \url{https://www.youtube.com/watch?v=BQ4yd2W50No}
\item \url{https://www.youtube.com/watch?v=StTqXEQ2l-Y}
\end{itemize}

\end{document}
